% =====================================================================
% CLASSP: Continuum Limit of Apollonian-Soddy Sphere Packing Dynamics
% Final version for ai.vixra upload --- May 2026
% =====================================================================
\documentclass[11pt,a4paper]{article}

\usepackage[margin=1in]{geometry}
\usepackage{amsmath,amssymb,amsfonts,amsthm}
\usepackage{hyperref}
\usepackage{booktabs}
\usepackage{array}
\usepackage{enumitem}
\usepackage{mdframed}

\hypersetup{
  colorlinks=true,
  linkcolor=blue,
  citecolor=blue,
  urlcolor=blue
}

\newtheorem{theorem}{Theorem}
\newtheorem{definition}{Definition}
\newtheorem{conjecture}{Conjecture}
\newtheorem{assumption}{Assumption}
\newtheorem{corollary}{Corollary}

\newcommand{\ASSP}{\mathrm{ASSP}}
\newcommand{\FTC}{\mathrm{FTC}}
\newcommand{\dd}{\mathrm{d}}
\newcommand{\kpc}{\mathrm{kpc}}
\newcommand{\kly}{\mathrm{kly}}
\newcommand{\R}{\mathbb{R}}

\title{CLASSP: The Continuum Limit of Apollonian--Soddy\\
Sphere Packing Dynamics as the Origin of Weak-Field\\
Gravity and Dirac Quantum Mechanics}
\author{Steven E.\ Elliott}
\date{May 2026}

\begin{document}
\maketitle

\begin{abstract}
The companion TFOFT paper proposes that the universe is a finite computational object defined by a dynamical Apollonian-Soddy sphere-packing fractal quine under Presburger arithmetic. This paper develops the first continuum-limit equations of that framework. The sphere radius is treated as a local clock field; its coarse-grained graph dynamics yields Newtonian gravity and the weak-field time component of general relativity. On the same tangency graph, a one-dimensional chiral transfer channel with a two-state boundary-spin register yields the 1+1 dimensional Dirac equation, with mass interpreted as the chirality-reversal rate per local radius-clock tick. The extension to full curved 3+1 dimensional spinor dynamics is stated as a frame-rotation ansatz, not a theorem. Electromagnetism is sketched as lower-layer callback momentum bookkeeping whose graph form is a candidate edge-phase holonomy limit. The paper also derives a corrected Milky Way prediction: if the Fermi bubbles are the polar lobes of a scaled 3dz2 boundary state, then a weaker equatorial "Fermi donut" should appear at roughly 10 to 15 kpc, with an inner shoulder near 8 to 12 kpc and one-quarter the angular density contrast of the polar lobes. This places the key observational test in the Gaia DR3 outer-disk regime.

\end{abstract}
\tableofcontents
\newpage

% ===================================================================
\section{Motivation}
% ===================================================================

The original TFOFT paper~\cite{TFOFT2026} proposes MATHICCS: every
mathematical operation in a physical law must correspond to a persistently
realizable physical procedure. Completed infinities, exact real continua, and
singular points are therefore not ontological primitives. They may remain
useful approximations, but only as coarse-grained descriptions of finite
physical computation.

The candidate finite ontology is the dynamical Apollonian--Soddy Sphere
Packing~\cite{Soddy1936,Boyd1973,Graham2003}. Primitive objects are not fields
on a continuum, but spheres, curvatures, tangencies, and local update rules. A
single sphere stores a state; two tangent spheres form the minimal
law-generating unit, the Gyro Ball. Across fractal scales, the same object is
conjectured to appear as electron/nucleus at atomic scale and star/core at
galactic scale. Flat galactic rotation curves~\cite{Sofue2012,Chemin2009} and
atomic orbital quantization are then interpreted as opposite-side projections
of one fractal boundary dynamics.

This paper asks a narrow technical question:
\begin{quote}
If the ASSP ontology of~\cite{TFOFT2026} is correct, what continuum field
equations appear when a bounded-resolution observer coarse-grains the finite
tangency dynamics?
\end{quote}
The answer developed here is that the two dominant continuum theories emerge
from complementary limits of the same graph:
\begin{align*}
\text{radius-clock scalar }u=\ln(R/R_0)
  &\;\longrightarrow\; \text{Newtonian gravity and weak-field }g_{00},\\
\text{directed tangency walk with Gyro mixing}
  &\;\longrightarrow\; \text{Dirac quantum matter},\\
\text{lower-layer callback momentum and U(1) holonomies}
  &\;\longrightarrow\; \text{candidate Maxwell/gauge field limit}.
\end{align*}
The target is not to quantize general relativity. Under TFOFT, that target is
category-mistaken: GR and QM are already effective views of the same finite
fractal update, observed from different sides of the boundary.

% ===================================================================
\section{The ASSP Tangency Graph}
% ===================================================================

\subsection{Graph structure}

Let the sphere packing define a graph $G=(V,E)$, where each vertex $i\in V$
is a sphere with center $x_i\in\R^3$, radius $R_i>0$, and curvature
$b_i=1/R_i$. An edge is a tangency:
\begin{equation}
  (i,j)\in E \iff |x_i-x_j|=R_i+R_j.
  \label{eq:tangency}
\end{equation}
The edge vector and length are
\begin{equation}
  s_{ij}=x_j-x_i,\qquad \ell_{ij}=|s_{ij}|=R_i+R_j.
  \label{eq:edge}
\end{equation}
The packing is recursively constrained by the oriented three-dimensional
Descartes--Soddy curvature relation~\cite{Soddy1936,Graham2003}.  For five
mutually tangent oriented spheres with signed curvatures $k_i=\pm R_i^{-1}$,
where an enclosing sphere carries negative curvature,
\begin{equation}
  \left(\sum_{i=1}^{5}k_i\right)^2
  =3\sum_{i=1}^{5}k_i^2.
  \label{eq:soddy_3d}
\end{equation}
Solving for a fifth curvature given $k_1,\ldots,k_4$, define
\begin{equation}
  S=\sum_{i=1}^{4}k_i,\qquad Q=\sum_{i=1}^{4}k_i^2.
\end{equation}
Then
\begin{equation}
  k_5=\frac{S\pm\sqrt{3S^2-6Q}}{2}.
  \label{eq:soddy_root}
\end{equation}
Equation~\eqref{eq:soddy_root} gives the two possible fifth spheres completing
one Descartes configuration.  Recursive Apollonian generation is more naturally
written as a \emph{reflection} rule on an existing five-sphere Descartes
configuration.  Replacing sphere $i$ while holding the other four fixed gives
\begin{equation}
  k_i'=\sum_{j\ne i}k_j-k_i.
  \label{eq:soddy_reflect_curv}
\end{equation}
The same reflection applies to curvature-weighted centers.  Writing
$m_i^a=k_i x_i^a$ for $a=1,2,3$,
\begin{equation}
  (m_i^a)'=\sum_{j\ne i}m_j^a-m_i^a,
  \qquad
  (x_i^a)'=\frac{(m_i^a)'}{k_i'}.
  \label{eq:soddy_reflect_center}
\end{equation}
Equations~\eqref{eq:soddy_reflect_curv}--\eqref{eq:soddy_reflect_center} are
the microscopic update rules used by the oriented 3D ASSP generator.  For a
tetrahedral seed of four unit spheres, the virtual enclosing sphere has
negative curvature
\begin{equation}
  k_{\rm out}=\frac{4-\sqrt{24}}{2}<0,
  \label{eq:outer_curvature}
\end{equation}
and the first reflected interior sphere has radius
\begin{equation}
  R_{\rm in}=\frac{1}{2+\sqrt{6}}\approx0.224745.
\end{equation}
The explicit curvature law fixes the microscopic recursive generation of the
packing.  The continuum-limit arguments below do not require this recursion
directly; they require the induced tangency graph, finite local second moments,
and statistical balance/isotropy after coarse-graining.

\subsection{Radius as local time}

The defining physical identification is
\begin{equation}
  u_i=\ln\frac{R_i}{R_0},
  \label{eq:u_def}
\end{equation}
where $R_0$ is the reference radius. In the Presburger-minimal TFOFT basis one
takes $R_0=\pi$~\cite{TFOFT2026}, but the continuum-limit derivations below do
not depend on that choice. The local clock rate is
\begin{equation}
  \dd\tau_i=e^{u_i}\dd t=\frac{R_i}{R_0}\dd t.
  \label{eq:clock}
\end{equation}
A bounded-resolution observer describes this through the emergent metric
\begin{equation}
  \dd s^2=-c^2e^{2u(x)}\dd t^2+\dd\ell_G^2,
  \label{eq:metric}
\end{equation}
where $\dd\ell_G^2$ is the spatial metric induced by the packing. In the GR
weak field, $\dd\tau\approx(1+\Phi/c^2)\dd t$; comparing with
$\dd\tau=e^u\dd t\approx(1+u)\dd t$ gives
\begin{equation}
  \Phi=c^2u=c^2\ln\frac{R}{R_0}.
  \label{eq:phi_u}
\end{equation}
\emph{Gravity is the coarse-grained logarithmic gradient of the ASSP
radius-clock field.}

\subsection{Normalized graph Laplacian}

For a field $f_i=f(x_i)$ slowly varying across the graph, define the weighted
graph Laplacian
\begin{equation}
  (\Delta_G f)_i=A_i\sum_{j\sim i}w_{ij}(f_j-f_i),
  \label{eq:lap_general}
\end{equation}
with symmetric tangency conductances $w_{ij}=w_{ji}>0$. A natural
radius-symmetric choice is
\begin{equation}
  w_{ij}=\frac{4\pi}{R_i+R_j}\left(\frac{R_iR_j}{R_i+R_j}\right)^2,
  \label{eq:weights}
\end{equation}
though the continuum proof below requires only local symmetry and finite
second moments. Define
\begin{equation}
  M_i^{ab}=\sum_{j\sim i}w_{ij}s_{ij}^a s_{ij}^b,\qquad
  M_i=\sum_{j\sim i}w_{ij}|s_{ij}|^2.
  \label{eq:moments}
\end{equation}
We assume local coarse-grained balance and isotropy:
\begin{align}
  \sum_{j\sim i}w_{ij}s_{ij}^a&=0,
  \label{eq:balance}\\
  M_i^{ab}&=\frac{M_i}{3}\delta^{ab}.
  \label{eq:isotropy}
\end{align}
Condition~\eqref{eq:balance} holds in idealized ASSP interiors by packing
symmetry; condition~\eqref{eq:isotropy} holds statistically under local
coarse-graining over an ASSP cluster~\cite{Mandelbrot1983,Boyd1973}. Choose
$A_i=6/M_i$, giving
\begin{equation}
  (\Delta_G f)_i
  =\frac{6}{\displaystyle\sum_{j\sim i}w_{ij}|s_{ij}|^2}
    \sum_{j\sim i}w_{ij}(f_j-f_i).
  \label{eq:normalized_lap}
\end{equation}

% ===================================================================
\section{Weak-Field Gravity as the Radius-Clock Continuum Limit}
% ===================================================================

\subsection{Coarse-graining theorem}

\begin{theorem}[Laplacian limit]
\label{thm:laplacian}
Let $f_i=f(x_i)$ with $f$ smooth on scales large compared to the local sphere
spacing $h_i\sim R_i$. Under conditions~\eqref{eq:balance}--\eqref{eq:isotropy},
\begin{equation}
  (\Delta_G f)_i=\nabla^2f(x_i)+O(h_i|\nabla^3f|).
  \label{eq:lap_limit}
\end{equation}
\end{theorem}

\begin{proof}
Taylor-expand $f_j=f(x_i+s_{ij})$:
\begin{equation}
  f_j-f_i=s_{ij}^a\partial_a f+\frac12s_{ij}^as_{ij}^b\partial_a\partial_bf
  +O(|s_{ij}|^3\nabla^3f).
  \label{eq:taylor}
\end{equation}
Substitute into~\eqref{eq:lap_general}. The first-order term vanishes
by~\eqref{eq:balance}. The leading contribution is
\begin{align}
  (\Delta_G f)_i
  &=\frac{A_i}{2}\sum_{j\sim i}w_{ij}s_{ij}^as_{ij}^b\partial_a\partial_bf
     +O(h_i|\nabla^3f|) \nonumber\\
  &=\frac{A_i}{2}\cdot\frac{M_i}{3}\delta^{ab}\partial_a\partial_bf
     +O(h_i|\nabla^3f|). \nonumber
\end{align}
Using $A_i=6/M_i$ gives $(\Delta_G f)_i=\nabla^2f+O(h_i|\nabla^3f|)$.
\end{proof}

\subsection{Poisson equation and Newtonian gravity}

Postulate the ASSP radius-clock source equation
\begin{equation}
  \Delta_Gu_i=\frac{4\pi G}{c^2}\rho_i.
  \label{eq:discrete_poisson}
\end{equation}
By Theorem~\ref{thm:laplacian}, the continuum limit is
$\nabla^2u=(4\pi G/c^2)\rho$. Multiplying by $c^2$ and using
$\Phi=c^2u$ from~\eqref{eq:phi_u}:
\begin{equation}
  \boxed{\nabla^2\Phi=4\pi G\rho.}
  \label{eq:poisson}
\end{equation}
The slow-particle geodesic acceleration follows from metric~\eqref{eq:metric}:
\begin{equation}
  \vec a=-\frac{c^2}{2}\nabla\ln\!\left(-\frac{g_{00}}{c^2}\right)
  =-c^2\nabla u=-\nabla\Phi.
  \label{eq:acceleration}
\end{equation}
For a spherical mass $M$, $\Phi(r)=-GM/r$, giving
\begin{equation}
  R(r)=R_0\exp\!\left(-\frac{GM}{c^2r}\right)
  \approx R_0\!\left(1-\frac{GM}{c^2r}\right),
  \label{eq:R_point}
\end{equation}
so local radius, local clock rate, and local coordinate light speed decrease
in a potential well.

\subsection{Weak-field GR slot}

From $\Phi=c^2u$ and metric~\eqref{eq:metric}:
\begin{equation}
  g_{00}=-c^2e^{2\Phi/c^2}
  \approx -c^2\!\left(1+\frac{2\Phi}{c^2}\right),
  \label{eq:g00_weak}
\end{equation}
which is the weak-field GR time component. The ASSP radius-clock scalar $u$
is the pre-continuum object whose smooth limit is the gravitational potential
and whose weak-field metric form is the $g_{00}$ sector of GR. The radius-clock
argument alone does not derive the full Einstein tensor; the tensor sector is
deferred to a later microscopic derivation based on spatial packing dynamics,
frame transport, and two-ball boundary exchange.


\subsection{Mass, inertia, and the deferred tensor sector}
\label{sec:no_eep_tensor_deferred}

The radius-clock calculation above recovers the Newtonian potential and the
weak-field time component of GR, but it does \emph{not} derive the full Einstein
tensor.  Earlier heuristic visual-relocation language is deliberately avoided
here because it introduces ambiguous displacement variables and stress-energy
normalizations before the microscopic tensor-sector dynamics has been fixed.

The CLASSP interpretation also does not assume the Einstein Equivalence
Principle as a primitive.  Instead:
\begin{align}
  \text{mass} &\equiv \text{a persistent angular-momentum register},
  \label{eq:mass_angular_register}\\
  \text{inertia} &\equiv \text{the linear-momentum cost required to translate
  that angular-momentum register}.
  \label{eq:inertia_linear_cost}
\end{align}
Thus gravitational response is not postulated as universal free fall in a
smooth spacetime.  It is the continuum appearance of the translation cost of
moving angular-momentum systems through a nonuniform radius-clock substrate.
The scalar sector is captured by
\begin{equation}
  \Phi=c^2\ln(R/R_0),\qquad \vec a=-\nabla\Phi,
\end{equation}
while the full tensor sector must come from the coarse-grained dynamics of
spatial packing geometry, frame transport, and boundary exchange.

The future tensor-sector task is therefore to derive an effective equation of
the schematic form
\begin{equation}
  G_{\mu\nu}^{\rm eff}
  =
  \frac{8\pi G}{c^4}T_{\mu\nu}^{\rm eff}
\end{equation}
from explicit ASSP update rules for curvature-weighted centers, tangency-frame
transport, and two-ball boundary exchange.  This paper does not claim that
result.  It claims only the controlled scalar result:
\begin{equation}
  \Delta_G u\to\nabla^2u,
  \qquad
  \Phi=c^2u,
  \qquad
  \nabla^2\Phi=4\pi G\rho.
\end{equation}

\subsection{M\"obius boundary acceleration and the Newton channel}

At a tangency point the M\"obius boundary map is conjectured to generate the
radial acceleration profile
\begin{equation}
  a_M(r)=\frac{2v_M^2}{r}\!\left(1+\frac{r_s}{r}\right)^{-2}
       =\frac{2v_M^2r}{(r+r_s)^2},\qquad
  r_s=\frac{2GM}{c^2},
  \label{eq:mobius_accel}
\end{equation}
where $v_M$ is the effective boundary-channel speed.  The Laurent expansion is
\begin{equation}
  a_M(r)=\frac{2v_M^2}{r}-\frac{4v_M^2r_s}{r^2}
  +\frac{6v_M^2r_s^2}{r^3}-\frac{8v_M^2r_s^3}{r^4}+O(r_s^4/r^5).
  \label{eq:laurent}
\end{equation}
The first correction term is
\begin{equation}
  a_{M,1}(r)=-\frac{4v_M^2r_s}{r^2}
  =-\frac{8v_M^2}{c^2}\frac{GM}{r^2}.
  \label{eq:mobius_newton_prefactor}
\end{equation}
Thus the M\"obius expansion has the Newtonian radial form with a dimensionless
prefactor $8v_M^2/c^2$.  In the FTC interpretation this prefactor is fixed by
the final map between the two tangent sphere surfaces.

Locally, each sphere surface is described by two coordinates $(x,y)$, and the
radial boundary coordinate needed by the tangency map is the Euclidean norm
\begin{equation}
  \rho=\sqrt{x^2+y^2}.
  \label{eq:surface_norm}
\end{equation}
A two-sphere tangency update therefore requires two such norm programs,
\begin{equation}
  \rho_1=\sqrt{x_1^2+y_1^2},\qquad
  \rho_2=\sqrt{x_2^2+y_2^2}.
\end{equation}
In the Presburger/FTC cost model, each local two-coordinate norm carries the
diagonal square-map cost $\sqrt2$.  The final two-surface boundary-map cost is
therefore
\begin{equation}
  C_{\rm map}=\sqrt2+\sqrt2=2\sqrt2.
  \label{eq:two_sqrt_cost}
\end{equation}
The normalized M\"obius boundary speed is the light-speed channel divided by
this final mapping cost:
\begin{equation}
  v_M=\frac{c}{C_{\rm map}}=\frac{c}{2\sqrt2},
  \qquad
  \frac{v_M^2}{c^2}=\frac{1}{8}.
  \label{eq:mobius_speed_norm}
\end{equation}
Substituting Eq.~\eqref{eq:mobius_speed_norm} into
Eq.~\eqref{eq:mobius_newton_prefactor} selects the Newtonian acceleration
normalization
\begin{equation}
  \boxed{a_{M,1}(r)=-\frac{GM}{r^2}.}
  \label{eq:newton_from_mobius}
\end{equation}
For a test mass $m$, this gives the Newton force law
\begin{equation}
  \boxed{F_N=-\frac{GMm}{r^2}.}
  \label{eq:newton_force}
\end{equation}
With the same normalization, the full expansion reads
\begin{equation}
  a_M(r)=\frac{c^2}{4r}-\frac{GM}{r^2}
  +\frac{3G^2M^2}{c^2r^3}
  -\frac{8G^3M^3}{c^4r^4}
  +O\!\left(\frac{G^4M^4}{c^6r^5}\right).
  \label{eq:mobius_normalized_expansion}
\end{equation}
The leading $1/r$ term is interpreted as the flat-rotation/inertial boundary
channel, while the first correction is the Newtonian channel selected by the
final two-surface norm-map cost.

\subsection{Dynamic radius field equation}

The dynamical ASSP radius equation is taken to be
\begin{equation}
  (\Delta_Gu)_i-\frac{1}{c^2}\ddot u_i=\frac{4\pi G}{c^2}\rho_i,
  \label{eq:wave_gravity}
\end{equation}
so that the static limit reproduces Eq.~\eqref{eq:discrete_poisson}.  With
\begin{equation}
  \Box=-\frac{1}{c^2}\partial_t^2+\nabla^2,
\end{equation}
the continuum limit is the scalar wave equation
\begin{equation}
  \Box u=\frac{4\pi G}{c^2}\rho.
  \label{eq:scalar_wave}
\end{equation}
Gravitational radiation corresponds, at this level, to propagating
perturbations of $u$. The tensorial sector is expected to arise from the
coarse-grained dynamics of $\dd\ell_G^2$.

% ===================================================================
\section{Quantum Matter as a One-Dimensional Tangency-Channel Limit}
% ===================================================================

\subsection{Why the Dirac channel is effectively one-dimensional}

The ASSP substrate is three-dimensional, but a single transfer event is not a
three-dimensional bulk motion.  It is a directed event along an active
tangency channel.  At the microscopic level an Info Ball either crosses the
local tangency boundary in one orientation or in the opposite orientation.
Thus the minimal quantum register is not a full spatial vector field; it is a
one-dimensional transport axis equipped with a two-state internal boundary
spin.

Let \(e\) denote the active tangency direction through a local pair of tangent
spheres.  The transfer coordinate \(x\) is distance along this channel.  The
two internal states are written
\[
  \psi_R(x,t),\qquad \psi_L(x,t),
\]
where \(R\) and \(L\) may be read equivalently as right/left movers,
clockwise/anti-clockwise boundary rotations, or the two signs of the local
spin register about the tangency hole.  The transverse ASSP geometry fixes the
local clock rate, frame orientation, and channel availability, but the local
propagation law itself is one-dimensional.

This is the key technical simplification: the Dirac structure does not first
require a complete derivation of a \(3+1\)-dimensional continuum spin bundle.
A single active tangency line plus a binary half-turn register is already
enough to generate the Dirac equation in the controlled continuum limit.

\subsection{Half-turn boundary register}

In TFOFT the rotational state channel of a Soddy sphere is hemisphere
restricted.  The rotational contribution to the fundamental state cost is
\[
  T_R=\pi r\big|_{r=\pi}=\pi^2,
\]
not \(2\pi^2\).  This is interpreted as a half-turn register: the physical
update need only distinguish the two boundary orientations of the active
rotational channel, not specify an unconstrained global rotation.

For CLASSP this means that the local quantum degree of freedom is naturally
binary.  Around a tangency hole the boundary spin decomposes into two
orientations,
\[
  \Psi =
  \begin{pmatrix}
    \psi_+\\
    \psi_-
  \end{pmatrix},
\]
which may also be written as \(\Psi=(\psi_R,\psi_L)^T\).  The two components
are not separate particles.  They are the two possible orientations of the same
one-dimensional transfer process.

\subsection{One-dimensional chiral walk}

Let \(\epsilon\) be the spatial step along the active tangency channel and
\(\Delta\tau\) the local proper tick.  The light-speed condition is
\[
  \frac{\epsilon}{\Delta\tau}=c.
\]
Let \(\theta\) be the local half-turn mixing angle.  The discrete update rule is
\begin{align}
  \psi_R(x+\epsilon,t+\Delta\tau)
    &=\cos\theta\,\psi_R(x,t)-i\sin\theta\,\psi_L(x,t),
  \label{eq:walk_R}\\
  \psi_L(x-\epsilon,t+\Delta\tau)
    &=\cos\theta\,\psi_L(x,t)-i\sin\theta\,\psi_R(x,t).
  \label{eq:walk_L}
\end{align}
Here the streaming part moves the two boundary orientations in opposite
directions along the channel, while the \(\theta\)-term flips the local
orientation.  In the FTC interpretation this flip is the local half-turn of
the boundary spin register.

For \(\theta\ll1\) and slowly varying amplitudes, Taylor expansion gives
\begin{align}
  (\partial_t+c\partial_x)\psi_R
    &=-i\frac{\theta}{\Delta\tau}\psi_L,
  \label{eq:chiral_R_theta}\\
  (\partial_t-c\partial_x)\psi_L
    &=-i\frac{\theta}{\Delta\tau}\psi_R.
  \label{eq:chiral_L_theta}
\end{align}
Identifying
\[
  \frac{\theta}{\Delta\tau}=\frac{mc^2}{\hbar}
\]
yields
\begin{align}
  (\partial_t+c\partial_x)\psi_R
    &=-i\frac{mc^2}{\hbar}\psi_L,
  \label{eq:chiral_R}\\
  (\partial_t-c\partial_x)\psi_L
    &=-i\frac{mc^2}{\hbar}\psi_R.
  \label{eq:chiral_L}
\end{align}
Writing \(\Psi=(\psi_R,\psi_L)^T\), this becomes
\[
  \boxed{
  i\hbar\partial_t\Psi
  =
  -i\hbar c\,\sigma_z\partial_x\Psi
  +mc^2\sigma_x\Psi.}
  \label{eq:dirac_1d}
\]
This is a representation-specific form of the \(1+1\)-dimensional Dirac
equation.  In CLASSP it is not merely a toy model: it is the continuum limit of
one active ASSP tangency-transfer channel.

\subsection{Radius-clock insertion}

On the ASSP graph the local proper tick is set by the local tangency step.  For
an edge \((i,j)\),
\[
  \Delta\tau_{ij}\sim\frac{R_i+R_j}{c}.
  \label{eq:tick_radius_edge}
\]
In a locally equal-radius coarse-grained region, \(R_j\approx R_i\), this
reduces to
\[
  \Delta\tau_i\sim\frac{2R_i}{c}.
  \label{eq:tick_radius}
\]
The mixing angle therefore becomes
\[
  \theta_i
  =
  \frac{mc^2\Delta\tau_i}{\hbar}
  \sim
  \frac{2mcR_i}{\hbar}
  =
  \frac{2R_i}{\lambda_C},
  \qquad
  \lambda_C=\frac{\hbar}{mc}.
  \label{eq:theta_radius}
\]
Thus mass is the chirality-reversal rate per local sphere-clock tick:
\[
  \boxed{
  m=\frac{\hbar}{c^2}\frac{\theta_i}{\Delta\tau_i}.}
  \label{eq:mass_flip}
\]
Massless modes have \(\theta=0\) and stream without local orientation reversal.
Massive modes repeatedly reverse the boundary-spin state, producing the
fractal zigzag path whose long-wavelength limit is Dirac-like.

\subsection{Radial-pole transfer and boundary spin}

For a near-radial transfer through a local pole of the packing, the relevant
coordinate is not the full three-dimensional displacement vector.  It is the
radial tangency coordinate selected by the local sphere pair.  The transfer
``cares'' about the boundary spin around the tangency hole: clockwise versus
anti-clockwise, or equivalently \(+\) versus \(-\).  Therefore the local quantum
question is binary:
\[
  \text{Which boundary orientation is active along this channel?}
\]
The answer is exactly the two-component spinor used above.

This is why the quantum update can be one-dimensional even though the packing
is three-dimensional.  The three-dimensional ASSP determines which local
channels exist and how their frames rotate from one sphere pair to the next,
but the elementary Dirac update only needs:
\begin{enumerate}
  \item one active propagation coordinate \(x\),
  \item one local clock tick \(\Delta\tau\),
  \item one binary half-turn spin register \((+,-)\).
\end{enumerate}

\subsection{Status of the three-dimensional extension}

  The result proved in this section is the effective one-dimensional channel
  limit:
  \[
    \text{active tangency line}
    +
    \text{two-state boundary spin}
    \quad\Longrightarrow\quad
    \text{\(1+1\)-dimensional Dirac equation}.
  \]
  The full \(3+1\)-dimensional curved-space Dirac equation is not derived here.
  It is a conjectural coarse-grained extension obtained by patching many local
  one-dimensional tangency channels together through the rotating ASSP frame.

  In that extension, a local tetrad \(e^\mu_a(x)\) records the averaged
  orientation of nearby tangency channels, and the spin connection records the
  frame change induced by passing from one sphere-pair channel to the next.  The
  expected continuum form is
  \[
    i\hbar\gamma^a e^\mu_a
    \left(\partial_\mu+\Omega_\mu\right)\Psi
    =
    mc\,\Psi,
    \label{eq:curved_dirac_ansatz}
  \]
  where \(\Omega_\mu\) is the effective spin connection generated by
  coarse-grained inter-sphere frame rotation.  Equation
  \eqref{eq:curved_dirac_ansatz} is therefore an ansatz, not a theorem.  The
  theorem-level result is the one-dimensional tangency-channel Dirac limit.

% ===================================================================
\section{Finite Tangency Registers, Bell Nonlocality, and Quantum Computation}
\label{sec:finite_quantum_registers}
% ===================================================================

\subsection{Superposition as finite tangency correlation}

The wavefunction is not taken here as an ontological primitive.  It is an
effective continuum bookkeeping object used by a bounded observer who cannot
resolve the full finite graph of mutual tangencies, chained tangencies,
cross-layer tangencies, and radius-clock ticks.

In the ASSP ontology, what ordinary quantum mechanics calls superposition is
interpreted as an unresolved finite set of compatible tangency configurations.
What ordinary quantum mechanics calls entanglement is interpreted as a linked
tangency-chain constraint.  Thus
\begin{equation}
  \boxed{
  \text{superposition}
  =
  \text{unresolved finite contact alternatives},
  }
  \label{eq:superposition_contact}
\end{equation}
and
\begin{equation}
  \boxed{
  \text{entanglement}
  =
  \text{linked tangency-chain constraint}.
  }
  \label{eq:entanglement_chain}
\end{equation}
Full local contact corresponds to maximal local correlation.  Lower-order
contact corresponds to lower-order correlation.  Long tangency chains
correspond to long-distance entanglement-like correlations, and cross-layer
tangencies correspond to apparent nonlocality relative to one emergent
spacetime layer.

There is no sharp quantum/classical boundary in this picture.  There is only a
change of coarse-graining.  At small register depth, finite tangency
alternatives must be treated explicitly and appear quantum-like.  At large
register depth, the same graph relations average into classical trajectories,
continuum fields, and effective probability densities.

\subsection{Bell correlations and cross-layer tangency}

CLASSP should not be read as a same-layer local hidden-variable theory.
Bell-type experiments rule out local hidden-variable explanations of quantum
correlations under the standard assumptions.  The ASSP proposal avoids that
category because the relevant hidden structure is not local with respect to
the emergent three-dimensional metric of one coarse-grained layer.

The underlying causal primitive is the directed tangency graph, including
same-layer tangency chains and cross-layer macrophoton/tangency channels.
Therefore two objects separated in emergent space may remain adjacent, linked,
or jointly constrained in the deeper ASSP graph.  The observed Bell correlation
is not produced by a signal moving through ordinary space after measurement.
It is produced by a prior or deeper graph relation whose causal adjacency is
not captured by the same-layer Lorentz cone.

Thus Bell's theorem is not violated.  Rather, one of the assumptions used to
derive Bell inequalities is not satisfied: locality with respect to the
emergent metric of one layer.  The model is explicitly nonlocal relative to
emergent spacetime.  It therefore belongs, if viable, to the class of nonlocal
hidden-structure theories, not local hidden-variable theories:
\begin{equation}
  \boxed{
  \text{Bell excludes same-layer local hidden variables, not cross-layer
  tangency adjacency.}
  }
  \label{eq:bell_cross_layer}
\end{equation}

No usable faster-than-light messaging is implied inside one layer, because
ordinary observers can only access the coarse-grained projection.  The deeper
graph may contain superluminal tangency adjacency, but controllable signaling
is filtered by register availability, channel orientation, radius-clock
mismatch, and the directed update order of the quine.

\subsection{Soddy contact bounds and local quantum registers}

The local ASSP quantum register is finite because sphere rearrangements are
constrained by sphere-packing contact laws.  In three Euclidean dimensions, the
strict Descartes--Soddy/Gosset closure contains five pairwise mutually tangent
oriented spheres:
\begin{equation}
  N_{\rm mutual}^{(3D)}=5.
  \label{eq:five_mutual_spheres}
\end{equation}
A distinct but physically important closure is the Soddy hexlet.  Given three
mutually tangent base spheres, a chain of six additional spheres closes around
them, with each chain sphere tangent to the three base spheres and to its two
chain-neighbors.  This gives a natural nine-sphere local contact register,
\begin{equation}
  N_{\rm hexlet}=3+6=9,
  \label{eq:hexlet_nine}
\end{equation}
without claiming that all nine spheres are pairwise mutually tangent.

In CLASSP, the five-sphere Descartes closure is interpreted as the maximum
local full-mutual-correlation register, while the nine-sphere hexlet is
interpreted as a finite local rearrangement neighborhood.  Tangency chains can
extend over arbitrary graph distance, but a single exact full-mutual local
correlation neighborhood is bounded.

The resulting hierarchy is
\begin{align}
  \text{pair tangency}
  &\Longrightarrow
  \text{two-register correlation}, \\
  \text{full mutual Descartes contact}
  &\Longrightarrow
  \text{maximal exact local correlation, capped at five registers}, \\
  \text{Soddy-hexlet closure}
  &\Longrightarrow
  \text{finite nine-sphere rearrangement neighborhood}, \\
  \text{long tangency chain}
  &\Longrightarrow
  \text{long-distance entanglement}, \\
  \text{cross-layer chain}
  &\Longrightarrow
  \text{apparent nonlocality beyond one-layer light cones}.
\end{align}

\subsection{Scale-relative causality and macrophoton transfer}
\label{sec:scale_relative_causality}

Special relativity is interpreted here as the effective causal geometry of one
coarse-grained fractal layer.  Its invariant speed \(c\) is the limiting
signal speed for ordinary bulk propagation inside that layer, measured using
that layer's radius-clock field.  It is not assumed to be the invariant speed
of every deeper or higher ASSP quine layer.

Let \(R_N\) denote the characteristic radius-clock scale of layer \(N\).  The
effective light speed measured inside that layer is written
\begin{equation}
  c_N\sim\frac{\Delta\ell_N}{\Delta\tau_N},
  \label{eq:layer_c}
\end{equation}
where both the distance unit \(\Delta\ell_N\) and the local tick
\(\Delta\tau_N\) are determined by the layer's ASSP geometry.  A transfer that
is ordinary and causal on layer \(N+1\) may project into layer \(N\) as an
apparently superluminal event:
\begin{equation}
  v_{\rm app}^{(N)}
  =
  \frac{\Delta x_N}{\Delta t_N}
  >
  c_N.
  \label{eq:apparent_superluminal}
\end{equation}
This does not imply time travel.  Causal order is not defined by the emergent
Lorentz cone of layer \(N\) alone, but by the directed update order of the
underlying tangency graph.  The graph update has a direction; the lower-layer
projection has an apparent distance and delay.  Apparent superluminality is
therefore a cross-layer projection effect, not a closed causal loop.

In the astrophysical case, a Star Ball forced into a Core Ball can excite a
macrophoton channel.  The corresponding emission need not reappear as a jet
from the same apparent galaxy.  In the ASSP graph, the receiving Core Ball may
be adjacent through a higher-layer tangency relation while being distant in
our emergent three-dimensional coordinates.  Thus an infall or callback event near one
Core-Ball register may correlate with a modulated jet return value at another
galaxy's core.  The process is nonlocal relative to the emergent metric of our
layer, but local relative to the directed tangency update of the higher layer.


\subsection{Core-Ball callback register and jet return values}
\label{sec:core_callback_register}

In CLASSP the object usually called a black-hole event horizon is not treated
as an ontologically primitive spacetime surface.  It is reinterpreted as a
\emph{Core-Ball callback register}: an addressable program boundary at which
remote Info-Ball transfers are received, decoded, and executed as local
precession, jet structure, and lower-layer memory allocation.  The programming
language is deliberate.  A local jet is not primarily local angular momentum
being expelled.  It is the visible return value of a remote angular-momentum
callback.

Let \(\Delta\vec L_{\rm rem}\) denote angular momentum carried by a mostly
one-dimensional Info-Ball/macrophoton transfer from a remote origin, and let
\(\hat q\) be the incoming transfer direction at the receiving Core Ball.  Let
\(\vec S_{\rm core}\) be the receiver spin register.  The callback torque is
modeled by the component of the remote instruction not already aligned with the
receiver spin:
\begin{equation}
  \vec\tau_{\rm cb}
  =
  \Gamma
  \left[
  \Delta\vec L_{\rm rem}
  -
  (\Delta\vec L_{\rm rem}\cdot\hat S_{\rm core})\hat S_{\rm core}
  \right],
  \label{eq:callback_torque}
\end{equation}
so that
\begin{equation}
  \frac{\dd\vec S_{\rm core}}{\dd t}=\vec\tau_{\rm cb}.
  \label{eq:core_spin_callback}
\end{equation}
The observed jet is the local visible trace of this receiver-side callback:
\begin{equation}
  \hat j_{\rm jet}\parallel\hat S_{\rm core}(t)
  \qquad\text{or}\qquad
  \hat j_{\rm jet}\parallel
  \frac{\dd\hat S_{\rm core}}{\dd t},
  \label{eq:jet_return_value}
\end{equation}
depending on whether the visible jet follows the instantaneous spin axis or
the precession axis.

The callback has an orientation-dependent routing efficiency.  If reflection is
most efficient along the polar spin axis, a minimal efficiency model is
\begin{equation}
  \eta_{\rm cb}^{\rm polar}=(\hat q\cdot\hat S_{\rm core})^2.
  \label{eq:callback_eff_polar}
\end{equation}
If reflection is most efficient through the receiver equatorial address plane,
then
\begin{equation}
  \eta_{\rm cb}^{\rm eq}=1-(\hat q\cdot\hat S_{\rm core})^2.
  \label{eq:callback_eff_equator}
\end{equation}
The core evolves to minimize callback routing cost,
\begin{equation}
  \mathcal C_{\rm cb}=1-\eta_{\rm cb},
  \label{eq:callback_cost}
\end{equation}
so spinning Core-Ball equators tend to self-align for efficient reflection of
remote Info-Ball transfers.

The channel split is the same rank-23 split used in TFOFT.  The visible return
value carries the addition-efficient channel,
\begin{equation}
  \Delta\vec L_{\rm return}=\frac{4}{27}\Delta\vec L_{\rm rem},
  \label{eq:return_channel_427}
\end{equation}
while the rank-23 overhead is written into the lower fractal layer as memory
allocation / inertial debt,
\begin{equation}
  \Delta\vec L_{\rm lower}=\frac{23}{27}\Delta\vec L_{\rm rem}.
  \label{eq:lower_channel_2327}
\end{equation}
Thus a jet is not exhaust from the local disk; it is a return value produced
when a receiving Core-Ball executes a remote angular-momentum callback.  The
associated dark/inertial structure is the lower-layer memory allocation needed
to execute that callback.

\subsection{Quantum Fourier Transform as a full-mutual-contact protocol}

Standard quantum computing treats the Hilbert space of \(n\) qubits as an
ideal \(2^n\)-dimensional complex vector space.  In CLASSP, this space is not
ontologically fundamental.  It is an effective bookkeeping representation of
finite contact and chained-contact relations in the underlying graph.

The Quantum Fourier Transform used in Shor-type order finding requires
coherent phase correlation across the active logical register.  In the ASSP
interpretation, exact unit-fidelity phase correlation is conjectured to require
full mutual tangency among the active logical qubit registers.  Lower-order
contact gives lower-order correlation; chained contact gives
entanglement-like correlation but not exact all-to-all phase coherence.

Because the strict three-dimensional Descartes--Soddy closure allows only five
pairwise mutually tangent oriented spheres, CLASSP makes the following
conjecture:
\begin{conjecture}[Five-register QFT coherence bound]
Exact full-mutual-contact Quantum Fourier Transform coherence can involve at
most five logical qubit-registers in a single local ASSP contact neighborhood:
\begin{equation}
  n_{\rm QFT}^{\rm exact}\leq 5.
  \label{eq:qft_five_bound}
\end{equation}
\end{conjecture}

For \(n\leq5\), a fully correlated QFT-like register may be physically
realizable in an ideal local contact neighborhood.  For \(n>5\), the
computation must use lower-order tangencies, chained tangencies, or routed
graph correlations.  These may still generate ordinary entanglement, but they
are not exact full-mutual-contact correlations.  The result is a predicted
structured error floor for large QFT-based algorithms.

The difference of sphere radii makes the limitation stronger.  Even within a
nominal five-register closure, unequal radii imply unequal local ticks,
\begin{equation}
  \Delta\tau_i\sim\frac{2R_i}{c}.
\end{equation}
A phase update with angular frequency \(\omega\) therefore accumulates a
relative radius-clock drift
\begin{equation}
  \delta\phi_{ij}
  \sim
  \omega(\Delta\tau_i-\Delta\tau_j)
  \sim
  \frac{2\omega}{c}(R_i-R_j).
  \label{eq:radius_phase_error}
\end{equation}
Thus full mutual tangency is necessary but not sufficient for perfect QFT
coherence.  Equal-radius or radius-compensated contact is also required.  In
real hardware this appears as structured phase error, not merely random
environmental decoherence.

\subsection{Shor limitation conjecture}

Shor's algorithm is mathematically valid in ideal Hilbert-space quantum
computation.  CLASSP does not dispute the formal algorithm.  It disputes the
assumption that arbitrary-size ideal QFT coherence is physically realizable as
an unbounded continuum Hilbert-space object.

\begin{conjecture}[CLASSP Shor limitation]
RSA-scale Shor factorization cannot be physically realized if exact order
finding requires full-mutual-contact QFT coherence beyond the five-register
Descartes bound.
\label{conj:shor_limit}
\end{conjecture}

Small compiled demonstrations are not inconsistent with this conjecture.  A
five-qubit demonstration of factoring \(N=21\), for example, lies exactly at
the proposed full-mutual-contact ceiling.  In the CLASSP reading, such
demonstrations are not early evidence of unlimited scalability; they may
instead be evidence that present demonstrations have reached the maximum local
contact order available to exact QFT-style phase correlation.

This is a falsifiable conjecture, not an established theorem.  If future
quantum computers run uncompiled, fault-tolerant Shor order finding at large
logical register sizes with no structured finite-contact error floor, then
Conjecture~\ref{conj:shor_limit} is falsified.  If instead large QFT circuits
exhibit persistent, nonthermal, geometry-like phase errors that cannot be
removed by standard error correction, the conjecture gains support.

\subsection{Experimental quantum-computing signatures}

Possible finite-register deviations from ideal quantum computation include:
\begin{enumerate}[label=\textbf{Q\arabic*.}]
  \item \textbf{Finite-register saturation:} large entangled states encounter
        nonrandom error floors associated with finite local contact closure.
  \item \textbf{Graph-topological decoherence:} errors correlate with hidden
        tangency-chain rearrangements rather than independent local noise.
  \item \textbf{Radius-clock phase noise:} chirality-flip phases drift with
        local radius-clock fluctuations.
  \item \textbf{Layer leakage:} strongly driven quantum systems leak
        correlation into adjacent fractal layers, appearing as anomalous
        decoherence or rare nonlocal synchronization.
  \item \textbf{QFT failure mode:} QFT-based algorithms show a transition from
        full-contact exact correlation to routed-contact approximate
        correlation once more than five logical registers require simultaneous
        all-to-all phase coherence.
\end{enumerate}
If quantum computation scales perfectly according to ideal Hilbert-space theory
with no finite-register anomalies, CLASSP loses support.  If large entangled
systems show structured, geometry-like, nonthermal, or nonlocal error
correlations, CLASSP gains a possible experimental foothold.

% ===================================================================
\section{Candidate Electromagnetic Limit from Edge-Phase Holonomies}
\label{sec:maxwell}
% ===================================================================


\subsection{Electromagnetism as lower-layer callback bookkeeping}
\label{sec:em_callback_bookkeeping}

The holonomy construction below gives the graph-theoretic form of the gauge
field.  The physical interpretation in TFOFT is that the vector potential is
the unresolved lower-layer momentum-flow register created by Core-Ball
callbacks.  Visible Info-Ball transfer carries the addition-efficient return
value, while the rank-23 overhead is written into the lower fractal layer as
orthogonal linear-momentum debt.

Let \(\Delta\vec L_{\rm rem}\) be the remote angular-momentum instruction
received at a Core-Ball callback register.  At callback radius \(r_{\rm cb}\),
with incoming direction \(\hat q\) and local radial direction \(\hat r\), define
an orthogonal lower-layer injection direction
\begin{equation}
  \hat n_\perp=
  \frac{\hat r\times\hat q}{|\hat r\times\hat q|}.
  \label{eq:orthogonal_injection_direction}
\end{equation}
The corresponding lower-layer linear-momentum write is modeled as
\begin{equation}
  \Delta\vec p_{\rm lower}
  =
  \frac{23}{27}
  \frac{|\Delta\vec L_{\rm rem}|}{r_{\rm cb}}
  \hat n_\perp.
  \label{eq:lower_momentum_write}
\end{equation}
Equivalently, relative to the visible return channel,
\begin{equation}
  \Delta\vec p_{\rm lower}
  =
  \frac{23}{4}
  \frac{|\Delta\vec L_{\rm return}|}{r_{\rm cb}}
  \hat n_\perp.
  \label{eq:lower_momentum_write_visible}
\end{equation}
This is not a local jet expelling local angular momentum.  It is the memory
write cost of executing a remote callback.

Let \(\vec\Pi_{\rm lower}\) be the coarse-grained lower-layer momentum density
produced by many such callbacks.  Define the effective FTC vector potential by
\begin{equation}
  \vec A_{\rm FTC}=\lambda_A\vec\Pi_{\rm lower}.
  \label{eq:A_FTC_lower_momentum}
\end{equation}
Then the magnetic-looking field is the curl of the lower-layer momentum-flow
register,
\begin{equation}
  \vec B_{\rm eff}=\nabla\times\vec A_{\rm FTC},
  \label{eq:B_curl_lower_momentum}
\end{equation}
and the electric-looking field is the time-changing callback pressure,
\begin{equation}
  \vec E_{\rm eff}=-\nabla\phi-\partial_t\vec A_{\rm FTC}.
  \label{eq:E_callback_pressure}
\end{equation}
Thus ordinary electromagnetic language is interpreted as continuum bookkeeping
for hidden lower-layer momentum allocation.  Intergalactic magnetic-field-line
behavior, jet alignment, and hypervelocity-star escape corridors are expected
to follow the same callback routing network rather than independent local
exhaust processes.

With the Lorenz-gauge condition
\begin{equation}
  \nabla\cdot\vec A_{\rm FTC}+\frac{1}{c^2}\partial_t\phi=0,
  \label{eq:lorenz_ftc}
\end{equation}
conservation of lower-layer callback momentum propagating at speed \(c\) gives
the candidate wave equation
\begin{equation}
  \left(\nabla^2-\frac{1}{c^2}\partial_t^2\right)\vec A_{\rm FTC}
  =-\mu_0\vec J_{\rm eff},
  \label{eq:A_wave_callback}
\end{equation}
where \(\vec J_{\rm eff}\) is the effective current associated with callback
momentum writes.  This is the physical route from remote angular-momentum
callbacks to the Maxwell holonomy limit.

\subsection{U(1) edge phases and discrete field strength}

Assign a U(1) phase to each oriented tangency edge.  Let $A_{ij}$ denote the
line integral of the continuum gauge potential along the edge,
\begin{equation}
  A_{ij}=\int_i^j A_\mu\,dx^\mu,\qquad A_{ij}=-A_{ji}.
\end{equation}
The edge phase is
\begin{equation}
  U_{ij}^{\rm EM}=\exp\!\left(\frac{iqA_{ij}}{\hbar}\right).
  \label{eq:edge_phase}
\end{equation}
The gauge-covariant graph derivative is
\begin{equation}
  (D_Af)_{ij}=U_{ij}^{\rm EM}f_j-f_i.
  \label{eq:cov_diff}
\end{equation}
The discrete flux phase is the holonomy around an elementary graph cycle
$C=(i_1,i_2,\ldots,i_k,i_1)$:
\begin{equation}
  F_C=\arg\!\prod_{(ij)\in C}U_{ij}^{\rm EM}
  =\frac{q}{\hbar}\sum_{(ij)\in C}A_{ij}.
  \label{eq:holonomy}
\end{equation}
For a minimal three-tangency cycle, $F_C$ is the dimensionless discrete analogue of magnetic flux through the triangle.  The corresponding continuum field-strength component is obtained only after dividing the unscaled line-integral sum by the oriented plaquette area in the coarse-grained limit.

\subsection{Discrete gauge action and Maxwell limit}

The candidate gauge-field action on the ASSP graph is
\begin{equation}
  S_A=-\frac14\sum_C W_C F_C^2+\sum_{(ij)}J_{ij}A_{ij},
  \label{eq:gauge_action}
\end{equation}
where $W_C$ are cycle weights set by the packing geometry. Stationarity
$\delta S_A/\delta A_{ij}=0$ gives a discrete divergence equation on the field
strength. The graph identity $d_G^2=0$ provides the discrete Bianchi identity
exactly. Under slow-variation and isotropy conditions, the cycle holonomies
expand to leading order in the gauge potential and the discrete divergence
maps to the continuum divergence, giving the candidate Maxwell limit
\begin{equation}
  \boxed{\partial_\mu F^{\mu\nu}=\mu_0J^\nu,\qquad
  \partial_{[\lambda}F_{\mu\nu]}=0.}
  \label{eq:maxwell}
\end{equation}
The normalization of $W_C$, the Lorentzian Hodge star, and the exact emergence
of $\mu_0$ are left as explicit tasks for the CLASSP gauge program.

\subsection{Coupling constant from register depth}

The electromagnetic coupling constant is conjectured to enter through the
Presburger register depth of one Soddy sphere~\cite{TFOFT2026,Soddy1936}:
\begin{equation}
  T_{\rm sphere}=4\pi^3+\pi^2+\pi\approx137.036304,
  \qquad
  \alpha\approx\frac{1}{T_{\rm sphere}}.
  \label{eq:alpha}
\end{equation}
Under the edge-phase ontology, $\alpha$ is the sequential-alignment probability
per quine cycle: the fraction of asynchronous two-ball register updates that
are causally ordered~\cite{TFOFT2026}. Higher-order tangency interactions are
conjectured to generate the gauge hierarchy toward $SU(2)$ and $SU(3)$.

% ===================================================================
\section{Simulation and Observation}
% ===================================================================

\subsection{Hydrogenic dark/inertial boundary fields}

In the TFOFT interpretation, dark matter is not primarily a particle species
but a lower-layer inertial boundary field~\cite{TFOFT2026}. The electron cloud
and the dark halo are both lower-layer boundary fields that constrain the
dynamics of the visible object (electron at atomic scale; stellar disk at
galactic scale). Therefore the simulation target is \emph{not}
\[
  \text{stars trace }|\psi_{nlm}|^2.
\]
The correct target is
\begin{equation}
  \mathcal{I}_{\rm halo}^{(nlm)}(R,z,\phi)
  \;\Longrightarrow\;
  \text{inertial/drag/clock field}
  \;\Longrightarrow\;
  \text{stellar disk response}.
  \label{eq:target}
\end{equation}
A minimal model uses hydrogenic orbital geometries as basis functions for the
dark/inertial field:
\begin{equation}
  \mathcal{I}(r,\theta,\phi)
  =\sum_{nlm}W_{nlm}
    \left|R_{nl}\!\left(\frac{r}{a_g}\right)Y_{lm}(\theta,\phi)\right|^2,
  \label{eq:inertia_basis}
\end{equation}
where $a_g$ is the galactic scale radius and $W_{nlm}$ are occupation weights.
A continuum observer may infer an effective density $\rho_{\rm eff}$ from
$\mathcal I$, but the physical object in the ASSP model is the inertial
boundary kernel rather than a literal cloud of particles.

The stellar disk evolves under
\begin{equation}
  \vec a_\star=-\nabla\Phi_{\rm baryon}-c^2\nabla u_{\rm halo}
  +\vec a_{\rm drag}[\mathcal I,\vec v]
  +\vec a_{\rm gyro}[\mathcal I,\vec v\times\hat n],
  \label{eq:disk_dynamics}
\end{equation}
where the first two terms are baryon and radius-clock gravity from
Eq.~\eqref{eq:poisson}, and the latter two represent the lower-layer inertial
drag and angular-momentum transfer channels responsible for flat-rotation
support. The M\"obius boundary acceleration~\eqref{eq:mobius_accel} provides
one analytic form for $\vec a_{\rm gyro}$.

For the $3d_{z^2}$ Milky Way state specifically, the effective inertial kernel
has the angular form
\begin{equation}
  \mathcal I_{3d_{z^2}}(r,\theta)
  =I_0\,\mathcal R_{32}(r/a_g)\,(3\cos^2\theta-1)^2,
  \label{eq:I_3dz2_general}
\end{equation}
where $\mathcal R_{32}$ is the scaled radial shell factor. If one uses the
hydrogenic density shape directly, then
\begin{equation}
  \mathcal R_{32}(r/a_g)
  \propto
  \left(\frac{r}{a_g}\right)^4
  \exp\!\left(-\frac{2r}{3a_g}\right).
  \label{eq:R32_shape}
\end{equation}
The scale $a_g$ is fixed by matching the polar Fermi-bubble height, not by
fitting the equatorial feature after the fact.  If $\mathcal I$ is treated as a
three-dimensional inertial density kernel, the radial density factor
$r^4e^{-2r/(3a_g)}$ peaks at $r=6a_g$.  If instead one fits the shell-weighted
radial measure $r^2\mathcal I$, the peak shifts to $r=9a_g$.  The simulations
must state which convention is used before fitting observational data.

\subsection{Numerical simulation program}

The simulation comparison program is:
\begin{enumerate}[label=\textbf{S\arabic*.}]
  \item Build disk simulations with four halo/inertial models: NFW, Einasto,
        generic dark disk, and the hydrogenic ASSP inertial kernel
        $\mathcal I_{3d_{z^2}}$ from Eq.~\eqref{eq:I_3dz2_general}.
  \item Fit rotation curves, vertical force $K_z$, HI gas flaring, spiral
        structure, and stream perturbations.
  \item Ask whether a small set of hydrogenic orbital boundary states explains
        the same observables with fewer effective parameters.
  \item Allow transitions between boundary states,
        $\Psi_{\rm halo}(t)=\sum_{nlm}a_{nlm}(t)\psi_{nlm}$, and identify the
        disk and plasma response to $\Delta|\Psi_{\rm halo}|^2$.
\end{enumerate}

\subsection{The $3d_{z^2}$ Milky Way state and the Fermi bubbles}

The companion paper~\cite{TFOFT2026} identifies the Fermi
bubbles~\cite{FermiBubbles2010,Ackermann2014} as the polar lobes of a scaled
$3d_{z^2}$ boundary state. The angular factor is
\begin{equation}
  Y_{20}(\theta)\propto3\cos^2\theta-1.
  \label{eq:Y20}
\end{equation}
Its nodal cones lie at
\begin{equation}
  3\cos^2\theta-1=0
  \;\Longrightarrow\;
  \theta_{\rm node}=\arccos\frac{1}{\sqrt3}\approx54.74^\circ.
  \label{eq:nodal_angle}
\end{equation}
The Fermi bubbles extend roughly tens of degrees above and below the Galactic
center; the common descriptive scale is about $50\,\kly\approx15.3\,\kpc$, while
some modeling papers quote heights of order $10\,\kpc$ depending on the
geometric convention and distance model~\cite{FermiBubbles2010,Ackermann2014}.
We therefore define
\begin{equation}
  H_{\rm FB}\equiv\text{polar Fermi-bubble height}\sim10\text{--}15\,\kpc.
  \label{eq:HFB}
\end{equation}
This scale correction is crucial: $50\,\kly$ is not $50\,\kpc$.

The TFOFT rank fraction $(23/27)$ may be read as setting the polar-lobe height
relative to the full local halo/boundary scale,
\begin{equation}
  H_{\rm FB}=\frac{23}{27}R_{\rm halo}^{\rm local}.
  \label{eq:HFB_rank}
\end{equation}
Thus, for $H_{\rm FB}\sim10$--$15\,\kpc$,
\begin{equation}
  R_{\rm halo}^{\rm local}\sim\frac{27}{23}H_{\rm FB}\sim12\text{--}18\,\kpc.
  \label{eq:Rhalo_local}
\end{equation}
This places the associated equatorial component in the outer visible disk,
not in a remote $40$--$50\,\kpc$ halo annulus.

\subsection{The equatorial ``Fermi donut'': corrected derivation and prediction}

At the poles and at the equator, the squared $Y_{20}$ amplitude gives:
\begin{align}
  |Y_{20}|^2_{\rm pole}&\propto(3\cos^2 0-1)^2=(2)^2=4,
  \label{eq:pole}\\
  |Y_{20}|^2_{\rm equator}&\propto(3\cos^2(\pi/2)-1)^2=(-1)^2=1.
  \label{eq:equator}
\end{align}
Therefore
\begin{equation}
  \boxed{\frac{\mathcal I_{\rm equator}}{\mathcal I_{\rm pole}}
  =\frac{|Y_{20}|^2_{\rm equator}}{|Y_{20}|^2_{\rm pole}}=\frac14.}
  \label{eq:quarter_density}
\end{equation}
The same $3d_{z^2}$ state that produces the polar Fermi bubbles requires an
equatorial toroidal lobe of one-quarter the polar angular density contrast.
In the simplest same-radius shell interpretation, the torus peak is at the
same orbital scale as the polar lobe:
\begin{equation}
  R_{\rm donut,peak}\sim H_{\rm FB}\sim10\text{--}15\,\kpc.
  \label{eq:donut_peak_corrected}
\end{equation}
The torus does not possess a unique hard inner edge because an orbital density
surface depends on the chosen contour. For a fiducial shell/shoulder estimate,
using the nodal-cone geometry gives an inner shoulder of order
\begin{equation}
  R_{\rm donut,in}\sim\sqrt{\frac{2}{3}}\,H_{\rm FB}
  \approx0.816H_{\rm FB}
  \sim8\text{--}12\,\kpc.
  \label{eq:donut_inner_corrected}
\end{equation}
This places the predicted feature from approximately the solar circle and
outer disk into the region just beyond the main stellar disk scale. The
feature is therefore not a far-halo object; it is an outer-disk boundary
signature.

\begin{mdframed}
\textbf{Corrected prediction: Fermi donut.}
Using the real Fermi-bubble scale $H_{\rm FB}\sim10$--$15\,\kpc$, the Milky
Way $3d_{z^2}$ boundary state predicts an equatorial annular counterpart with:
\begin{itemize}
  \item fiducial inner shoulder $R_{\rm in}\sim8$--$12\,\kpc$;
  \item radial peak $R_{\rm peak}\sim10$--$15\,\kpc$;
  \item angular density/inertia contrast $\approx1/4$ of the polar lobes;
  \item geometry: a broad, low-contrast ring symmetric about the Galactic plane.
\end{itemize}
A coherent outer-disk boundary at this scale supports the $3d_{z^2}$ model.
A smooth featureless disk/halo through $R\sim8$--$20\,\kpc$ disfavors this
particular state identification.
\end{mdframed}

\subsection{Directed disk-halo search cone and macro-photon leakage analogy}
\label{sec:directed_donut_search}

The annular counterpart is not predicted as an isotropic halo shell.  In the
present interpretation it is a directed disk-halo boundary associated with the
Sgr~A* Core-Ball channel.  The observational search should therefore emphasize
structures moving or aligned radially away from the Galactic center through the
outer disk, especially where the outward channel intersects the predicted
$R\sim8$--$20\,\kpc$ annulus.

The nodal-cone geometry also gives an angular search scale about the Galactic
plane.  The nodal angle measured from the polar axis is
\begin{equation}
  \theta_{\rm node}=\arccos\frac{1}{\sqrt3}\approx54.74^\circ.
\end{equation}
The corresponding elevation above the equatorial plane is
\begin{equation}
  \beta_{\rm node}
  =90^\circ-\theta_{\rm node}
  =\arcsin\frac{1}{\sqrt3}
  \approx35.26^\circ.
  \label{eq:equatorial_elevation}
\end{equation}
Thus the disk component should be searched as a broad low-contrast band with
angular support of order $34^\circ$--$35^\circ$ above and below the Sgr~A*
equatorial plane, not as a razor-thin disk.  Because the equatorial component
has only one-quarter the polar angular density contrast, the signal is expected
to be centrally weighted rather than uniform.  A simple working model is a
normal angular profile
\begin{equation}
  P(\beta)\propto
  \exp\!\left(-\frac{\beta^2}{2\sigma_\beta^2}\right),
  \qquad
  \sigma_\beta\approx12^\circ,
  \label{eq:donut_normal_angle}
\end{equation}
so that approximately $68$--$69\%$ of the equatorial leakage lies inside
$|\beta|\lesssim12^\circ$.

This provides a speculative analogy with ultra-high-energy ``Oh-My-God''-type
particle events.  If an event stream lies near $\beta\approx-14^\circ$
relative to the Sgr~A* equatorial plane, it is near a one-sigma excursion in
Eq.~\eqref{eq:donut_normal_angle}.  In the FTC interpretation such a stream
could be read as a small-scale Rutherford-gold-foil analogue: a macro-photon
or micro-macro-photon transfer from the Sgr~A* Core-Ball channel is deflected
by the lower-layer dark/inertial disk boundary.  The one-sided sign of the
observed offset would then encode the polarity of the emitted info-ball channel
relative to the lower-layer dark-matter leakage channel.

Equivalently, the redshift of a macrophoton is modeled as a momentum leak from
the initial ejection into the lower fractal substrate:
\begin{equation}
  \Delta p_{\rm macro\gamma}
  =
  -\Delta p_{\rm lower},
  \qquad
  \Delta p_{\rm lower}
  \parallel
  \nabla\mathcal I_{3d_{z^2}}.
  \label{eq:macrophoton_leak}
\end{equation}
This subsection is a proposed observational discriminator, not a detection
claim.  It predicts that candidate events should correlate with outward
Sgr~A*-centered trajectories, with the low-contrast equatorial dark boundary,
and with the sign/polarity of the Core-Ball to Info-Ball channel.

\subsection{Connection to Gaia DR3 outer-disk anomalies}

The corrected Fermi-donut scale overlaps the radial domain where recent Gaia
DR3 studies already report model-sensitive Milky Way dynamics. Sylos Labini
constructed generalized rotation curves from Gaia DR3 data over
$R=8.5$--$25\,\kpc$ and heights $z=-2$ to $+2\,\kpc$, comparing a canonical
NFW halo model with a dark-matter-disk (DMD) model in which the dark component
is confined to the Galactic plane and follows neutral hydrogen; that analysis
reported that the DMD model generally fits the large-radius on-plane and
off-plane behavior better than NFW~\cite{SylosLabini2024}. Other Gaia DR3
Jeans analyses find evidence for a declining rotation curve beyond
$\sim15\,\kpc$, while also warning that disequilibrium, non-axisymmetry, and
satellite perturbations can bias Jeans-inferred circular velocities at the
10--15\% level~\cite{Koop2024}. Earlier Gaia DR3 mapping extends disk
kinematic maps and rotation-curve reconstruction to roughly $30\,\kpc$,
providing exactly the radial coverage needed to search for the corrected
$3d_{z^2}$ equatorial boundary~\cite{Wang2023}.

This is not a claim that Gaia has already detected the Fermi donut. The point
is sharper: once the incorrect $50\,\kpc$ scale is removed, the TFOFT
prediction lands in the same $R\sim10$--$20\,\kpc$ outer-disk regime where
Gaia already indicates that simple spherical-halo modeling is under stress or
at least model-dependent. The $3d_{z^2}$ model predicts a specific angular and
radial boundary structure in that region, not merely a generic dark disk.

\subsection{Emission and star-formation scaling}

If the polar Fermi bubbles arise where the boundary state opens a high-density
angular channel, the equatorial torus produces weaker signatures for two
independent reasons: (1) the angular density is lower by
Eq.~\eqref{eq:quarter_density}; and (2) local baryonic plasma injection in the
outer disk differs from the central polar channel. If emissivity scales
linearly with boundary density and local star-formation/plasma production,
\begin{equation}
  I_{\rm donut}
  \sim\frac14
  \left(\frac{\Sigma_{\rm SFR}(R_{\rm donut})}{\Sigma_{\rm SFR,polar}}\right)
  I_{\rm bubble}.
  \label{eq:emission_linear}
\end{equation}
If the relevant emission is collisional or plasma-density squared,
\begin{equation}
  I_{\rm donut}
  \sim\frac{1}{16}
  \left(\frac{\Sigma_{\rm gas}(R_{\rm donut})}{\Sigma_{\rm gas,polar}}\right)^2
  I_{\rm bubble}.
  \label{eq:emission_quadratic}
\end{equation}
If the boundary operates by a threshold/ignition rule, the donut may appear
more strongly in dynamics, HI flaring, or stellar-stream perturbations than in
X-ray or gamma-ray brightness. Non-detection in all-sky emission maps is
therefore not decisive unless the search is tuned to the predicted low-contrast
outer-disk annulus.

\subsection{Observational program}

\begin{enumerate}[label=\textbf{O\arabic*.}]
  \item Search Gaia DR3-derived rotation and vertical-force residuals for a
        coherent boundary or slope change in $R\sim8$--$20\,\kpc$, especially
        near $R\sim10$--$15\,\kpc$, with the search geometry oriented outward
        from the Galactic center.
  \item Compare the $3d_{z^2}$ inertial kernel against NFW, Einasto, and
        generic dark-disk models using the same Gaia DR3 generalized rotation
        curves and off-plane constraints.
  \item Search HI and outer-disk gas maps for flaring, compression, or ringlike
        behavior at the corrected donut scale.
  \item Search Fermi/eROSITA/radio residual maps for weak annular nonthermal
        emission in the Galactic plane at the predicted scale.
  \item Test whether stellar streams, open clusters, globular clusters, or
        satellite orbits show perturbations when crossing the predicted
        annular boundary.
  \item Test whether ultra-high-energy particle or photon-like event directions
        cluster near the outward Sgr~A* equatorial channel, especially near
        one-sigma offsets such as $\beta\sim-12^\circ$ to $-14^\circ$.
\end{enumerate}

% ===================================================================
\section{Conclusion}
% ===================================================================

This paper gives the first continuum-limit skeleton of the TFOFT/ASSP ontology
of~\cite{TFOFT2026}. The primitive variable is the sphere radius; its logarithm
is the local clock field:
\begin{equation}
  u=\ln\frac{R}{R_0}.
\end{equation}
The radius-clock field obeys a graph Poisson equation that coarse-grains to
Newtonian gravity and to the weak-field $g_{00}$ sector of GR; independently,
the M\"obius boundary expansion recovers the Newton force coefficient when the
final two-surface norm-map cost is $C_{\rm map}=2\sqrt2$:
\begin{equation}
  \Delta_Gu=\frac{4\pi G}{c^2}\rho
  \;\Longrightarrow\;
  \nabla^2\Phi=4\pi G\rho,
  \qquad g_{00}=-c^2e^{2\Phi/c^2}.
\end{equation}
\begin{equation}
  C_{\rm map}=2\sqrt2,
  \qquad
  v_M=\frac{c}{2\sqrt2}
  \;\Longrightarrow\;
  a_{M,1}=-\frac{GM}{r^2}.
\end{equation}
On the same graph, directed light-speed tangency walks with Gyro-Ball
chirality mixing coarse-grain to the Dirac equation~\cite{Ord1983,Dirac1928},
with the spin connection geometrically induced by inter-sphere frame rotation.
Lower-layer callback momentum bookkeeping and U(1) edge-phase holonomies
provide a candidate graph route to Maxwell electromagnetism.  The tensor sector
is not derived here; it is deferred to a microscopic treatment of spatial
packing dynamics, frame transport, and two-ball boundary exchange.

GR and QM are not fundamental competitors in this ontology. They are effective
continuum shadows of a single finite fractal sphere-packing computation:
radius-clock gradients become gravity, directed tangency zigzags become
quantum matter, and lower-layer callback momentum/edge-phase windings become
gauge fields.

The immediate empirical test is galactic. Correcting the Fermi-bubble scale
from $50\,\kpc$ to $50\,\kly\approx15\,\kpc$ moves the predicted equatorial
$3d_{z^2}$ ``Fermi donut'' into the outer visible disk, with a fiducial inner
shoulder near $8$--$12\,\kpc$, a peak near $10$--$15\,\kpc$, one-quarter
the polar-lobe angular density contrast, and a preferred outward orientation
from the Galactic center. This is exactly the region where Gaia
DR3 rotation-curve and dark-disk analyses already show nontrivial,
model-sensitive behavior. The next step is to simulate spiral galaxies using
hydrogenic ASSP inertial kernels and to test the corrected $3d_{z^2}$ boundary
against Gaia, HI, stream, Fermi, eROSITA, and radio data before accepting the
full ontology.

\section*{Acknowledgments}
The author thanks the survey teams whose archival data will enable the
observational program of Section~6, and the mathematical physics community
whose work on Apollonian packings~\cite{Boyd1973,Graham2003}, graph
Laplacians, and Ord-style walks~\cite{Ord1983} provided the technical
scaffolding for the derivations above.

% ===================================================================
\begin{thebibliography}{99}

\bibitem{TFOFT2026}
S.~E. Elliott,
``The Fact of Fractal Tennis: The Universe as a Fractal Computer Defined by
the Star=Electron+Atom=Galaxy Quine,''
ai.vixra preprint 2605.0003, 2026.
\url{https://ai.vixra.org/abs/2605.0003}

\bibitem{Ord1983}
G.~N. Ord,
``Fractal space-time: a geometric analogue of relativistic quantum mechanics,''
\textit{J. Phys. A: Math. Gen.} \textbf{16}, 1869 (1983).

\bibitem{Soddy1936}
F.~Soddy,
``The Kiss Precise,''
\textit{Nature} \textbf{137}, 1021 (1936).

\bibitem{Boyd1973}
D.~W. Boyd,
``The residual set dimension of the Apollonian packing,''
\textit{Mathematika} \textbf{20}, 170 (1973).

\bibitem{Graham2003}
R.~L. Graham, J.~C. Lagarias, C.~L. Mallows, A.~R. Wilks, and C.~H. Yan,
``Apollonian circle packings: number theory,''
\textit{J. Number Theory} \textbf{100}, 1 (2003).

\bibitem{Mandelbrot1983}
B.~B. Mandelbrot,
\textit{The Fractal Geometry of Nature}
(W.~H. Freeman, 1983).

\bibitem{Presburger1929}
M.~Presburger,
``\"Uber die Vollst\"andigkeit eines gewissen Systems der Arithmetik ganzer
Zahlen,''
\textit{Comptes Rendus du I Congr\`es des Math\'ematiciens des Pays Slaves},
Warsaw, 1929, pp.~92--101.

\bibitem{Dirac1928}
P.~A.~M. Dirac,
``The quantum theory of the electron,''
\textit{Proc. R. Soc. A} \textbf{117}, 610 (1928).

\bibitem{FermiBubbles2010}
M.~Su, T.~R. Slatyer, and D.~P. Finkbeiner,
``Giant Gamma-ray Bubbles from Fermi-LAT: AGN Activity or Bipolar Galactic Wind?''
\textit{ApJ} \textbf{724}, 1044 (2010).

\bibitem{Ackermann2014}
M.~Ackermann et al.\ (Fermi-LAT Collaboration),
``The Spectrum and Morphology of the Fermi Bubbles,''
\textit{ApJ} \textbf{793}, 64 (2014).

\bibitem{Sarkar2015}
K.~C. Sarkar, B.~Nath, and P.~Sharma,
``On the origin of the Fermi bubbles,''
\textit{MNRAS} \textbf{453}, 3827 (2015).

\bibitem{Sofue2012}
Y.~Sofue,
``Rotation Curve and Mass Distribution in the Galactic Center,''
\textit{PASJ} \textbf{64}(4), 75 (2012).

\bibitem{Chemin2009}
L.~Chemin, C.~Carignan, and T.~Foster,
``Rotation Curve of the Andromeda Galaxy from Combined 21-cm Data,''
\textit{ApJ} \textbf{705}(2), 1395 (2009).

\bibitem{SylosLabini2024}
F.~Sylos Labini,
``Generalized rotation curves of the Milky Way from the GAIA DR3 data-set:
constraints on mass models,''
\textit{Astrophysical Journal} \textbf{976}, 185 (2024); arXiv:2410.14307.

\bibitem{Koop2024}
O.~Koop, T.~Antoja, A.~Helmi, T.~M.~Callingham, and C.~F.~P.~Laporte,
``On the Galactic rotation curve inferred from the Jeans equations: assessing
its robustness using Gaia DR3 and cosmological simulations,''
arXiv:2405.19028 (2024).

\bibitem{Wang2023}
H.-F. Wang, Z. Chrob\'akov\'a, M. L\'opez-Corredoira, and F. Sylos Labini,
``Mapping the Milky Way Disk with Gaia DR3: 3D Extended Kinematic Maps and
Rotation Curve to $\approx30$ kpc,''
\textit{Astrophysical Journal} \textbf{942}, 12 (2023).

\bibitem{Strassen1969}
V.~Strassen,
``Gaussian elimination is not optimal,''
\textit{Numer. Math.} \textbf{13}, 354 (1969).

\bibitem{CODATA2022}
E.~Tiesinga et al.,
``CODATA recommended values of the fundamental physical constants: 2022,''
\textit{Rev. Mod. Phys.} \textbf{96}, 025002 (2024).

\end{thebibliography}

\end{document}